\section{Flow matching and trajectory parameterization}
\label{sec:C06:flow}
Flow matching learns a time-dependent vector field that transports a simple distribution toward the data distribution~\cite{lipman2022flowmatching}. For a simple illustrative path $x_t=(1-t)x_0+tx_1$, with endpoints drawn from a chosen coupling, the conditional velocity is $x_1-x_0$. A regression objective trains $v_\theta(x_t,t)$ to predict that velocity. The marginally correct field is the conditional expectation of the target velocity given $(x_t,t)$, obtained by the usual squared-loss projection property. This formulation separates a training path from the numerical solver used at sampling time.

The interpolation variable in this construction is not physical environment time. It parameterizes a transformation between distributions. Confusing the two can lead to an incorrect interpretation of an intermediate generated state as a physically possible step in an agent trajectory. A model can generate an entire action sequence or future video as one sample while its internal denoising or transport path passes through objects with no environmental meaning. Planning constraints must apply to the generated trajectory's action and state coordinates, not to arbitrary intermediate solver iterates.

Numerical integration introduces another approximation. A learned vector field can be accurate under its training distribution while a coarse solver accumulates error or evaluates off-distribution states. More solver steps may improve fidelity but cost latency. The solver, tolerance, precision, and conditioning policy belong to the reproducible method. A claimed model improvement that changes both the learned field and the solver cannot isolate which component produced the gain. For agents, the complete outcome also depends on the external verifier that accepts or rejects the generated plan.

Action constraints can be enforced through parameterization, constrained sampling, projection, or post-generation rejection. Each changes the distribution or feasible set. Projection may alter a generated trajectory so that its predicted consequences no longer correspond to its actions. Rejection can require many expensive samples when feasible plans are rare. A parameterization can guarantee local bounds, such as a control range, while leaving global requirements unproven. These mechanisms should be tested against exact environment constraints where possible, rather than assuming that a high-quality generator learns every hard rule from examples.
